\section{Type Erasure}
The dynamic semantics given in Section \ref{sec:Formalism}
relies on ubiquitous type annotations
in order to determine how function application will proceed
or how a frame of sub-arrays should collapse to a single array.
While the possible case of constructing a frame
with no actual result cells whose shape can be inspected
can only be resolved by consulting a type annotation,
the types themselves contain more information than is strictly needed.
For example, it does not matter whether we are collapsing
an empty frame of functions,
an empty frame of integers,
or an empty frame of boxes.
The result shape is the same,
regardless of the type of the atoms contained within the cells.
All we truly need is the resulting shape (alternatively, the result cells' shape).
Similarly, evaluating a function application requires knowing
the expected cell shapes for the arguments,
but it could, in principle, be done
without knowing anything about their atoms.
Function application is still tagged with a result shape,
again to head off issues arising from mapping over an empty frame.

In a type-erased version of Remora,
we only need the term and index levels---%
the syntactic class of types is discarded.
The syntax for erased Remora is given in Figure \ref{fig:ErasedAbstractSyntax}.
Note that the grammar of type indices from Figure \ref{fig:AbstractSyntax} is still in use here,
although expressions, atoms, and their corresponding function and value-form subsets
are now replaced with type-erased versions.

\FigErasedAbstractSyntax

Evaluation in erased Remora proceeds similarly to explicit Remora.
A function-application form has a principal frame
chosen to be the largest of the function and argument frames,
and a $\mathit{lift}$ reduction replicates the function and argument arrays' atoms
to bring all of the frames into agreement.
The argument frames themselves are identified based on
the individual argument positions' cell-shape annotations,
rather than by inspecting a type annotation on the array in function position.
A $\mathit{map}$ reduction turns an application form where all pieces have the same frame
into a {\tt frame} form,
where the end-result shape matches the result shape tag on the original application.
Index application also maps over an array of index functions,
producing a {\tt frame} of substituted function bodies.
Since the type level has been eliminated,
there are no {\tt t$\lambda$} and {\tt t-app} forms
and no need for a $t\beta$ reduction rule.

\FigErasedSemantics

The translation from explicit Remora to erased Remora consists of three erasure functions:
$\EraseE{\cdot} : \mathit{Expr} \rightarrow \erased{\mathit{Expr}}$,
$\EraseA{\cdot} : \mathit{Atom} \rightarrow \erased{\mathit{Atom}}$,
and $\EraseT{\cdot} : \mathit{Type} \rightarrow \mathit{Index}$.
These functions are defined in Figure \ref{fig:ErasureDefn}.

\ErasureDefn

We also define $\EraseC{\cdot} : \Ctxt \rightarrow \erased{\Ctxt}$,
given in Figure \ref{fig:ContextErasure},
which is not needed for defining the erased form of an explicit Remora program
but is useful for demonstrating their equivalence.

\ContextErasure

Types in explicit Remora are turned into indices in erased Remora.
These indices are the dynamic residue of types,
in the same sense that term-level values are dynamic,
though the are still subject to a static discipline which governs their values
and their relation to the array values they describe.
Array types become just the shapes used to construct them,
whereas functions, universals, dependent sums and products, and base types
become the ``scalar'' shape.
Extracting the index components of all types means that type variables
can be turned into index variables,
which will stand for the index component of whatever type the variable originally stood for.
This translation captures exactly the information that a {\tt frame} form needs
in the event that there are no cells.
By extension, the term and index application forms
also get the bookkeeping information needed by
the {\tt frame}s they will eventually become.
% The non-lifting forms {\tt array} and {\tt unbox} can safely ignore the source term's type tag.

For example, consider a function term whose type is
{\tt (-> (s (Arr t (Shp))) (Arr t (Shp k)))},
where {\tt s}, {\tt t}, and {\tt k} are
bound as {\tt Array}, {\tt Atom}, and {\tt Dim} respectively.
This function produces a vector of some statically uncertain length
containing atoms of uncertain type.
When we apply this function,
the explicitly typed application form
describes the resulting array's type.
If our arguments are a single {\tt s}
and a ${\tt n} \times 4$ matrix of numbers,
with {\tt n} also bound as a {\tt Dim},
the principal frame shape is {\tt (Shp n 4)}.
So we will have result type {\tt (Arr Num (Shp n 4 k))}.
Type-erasing the application form must still preserve
enough information to produce an array of the correct shape,
even if {\tt n} turns out to be 0,
leaving us with no result cells whose shape we can inspect.
However, the dynamic semantics does not rely on knowing that
the result array contains {\tt Num}s.
% When type erasure recurs down the function position,
% we will effectively summarize the function array's type as ``scalar,''
% throwing out any information about the function's shape-level behavior.
% That information is only used for the whole application form's result,
% which we give as {\tt (Shp n 4 k)}.
The binders for index variables {\tt n} and {\tt k},
which must be either {\tt I$\lambda$} or {\tt unbox},
are still present in the type-erased program,
since the indices they eventually bind to those variables
will affect the program's semantics.
The {\tt T$\lambda$}s which bind {\tt s} and {\tt t}
turn into {\tt I$\lambda$}s,
though the variable {\tt t} is never used in the type-erased program.
If any type argument was bound to {\tt s} in the original program,
we replace it with its shape.
All occurrences of {\tt s} from the original program
now stand for an array shape rather than a full array type.

We develop a bisimulation argument to show that
the behavior of an explicitly typed term matches the behavior of its erased form.
We define the space $S$ of machine states to be the sum of
the set of well-typed explicit Remora terms
and the set of their type-erased forms.
That is, $S = \wt{Expr} \uplus \erased{\wt{\Expr}}$,
where
$\wt{Expr} =
\{\expr \in \Expr | \typeof{\emptyEnv}{\emptyEnv}{\emptyEnv}{\expr}{\type}\}$
and
$\erased{\wt{Expr}} =
\{\EraseE{\expr} | \expr \in \wt{\Expr}\}$.
Transitions in the machine match the explicit and erased languages' respective $\mapsto$ relations.
We also define the ``erasure equivalence'' relation $\EraseEqv$ on machine states as
the equivalence closure of the relation imposed by $\EraseE{\cdot}$.
Before we show that $\EraseEqv$ is a bisimulation, several intermediate results are needed.

First, the bisimulation proof will in one case need to
reach deep into an expression to find the next redex.
A compositionality property of the erasure rule
will make it possible to reason about the redex and its reduced form
separately from the evaluation context in which it is embedded.

\begin{lemma}[Erasure in context]
  \label{ErasureInContext}
  Given an evaluation context $\ctxt$ and expression $\expr$,
%  where $\typeof{\emptyEnv}{\emptyEnv}{\emptyEnv}{\fillc{\ctxt}{\expr}}{\type}$,
  where ${\fillc{\ctxt}{\expr}}$ is well-typed,
  $\EraseE{\fillc{\ctxt}{\expr}} = \fillc{\EraseC{\ctxt}}{\EraseE{\expr}}$.
\end{lemma}
\begin{sproof}[This follows from straightforward induction on $\ctxt$.]
  We use induction on $\ctxt$.

\paragraph{Array literal containing unevaluated box:}
  \[\ctxt = \arrlit{\sequence{\atval}\,\dsum{\sequence{\idx}}{\ctxt'}{\type}\,\sequence{\atom}}{\sequence{\nat}}\]
  $\EraseE{\fillc
    {\arrlit{\sequence{\atval}\,\dsum{\sequence{\idx}}{\ctxt'}{\type}\,\sequence{\atom}}{\sequence{\nat}}}
    {\expr}}$ \\
  $= \EraseE{\arrlit
    {\sequence{\atval}\,\dsum{\sequence{\idx}}{\fillc{\ctxt'}{\expr}}{\type}\,\sequence{\atom}}
    {\sequence{\nat}}}$ \\
  $= \arrlit
  {\sequence{\EraseA{\atval}}\,\EraseA{\dsum{\sequence{\idx}}{\fillc{\ctxt'}{\expr}}{\type}}\,\sequence{\EraseA{\atom}}}
  {\sequence{\nat}}$ \\
  $= \arrlit
  {\sequence{\EraseA{\atval}}\,\dsum{\sequence{\idx}}{\EraseE{\fillc{\ctxt'}{\expr}}}{}\,\sequence{\EraseA{\atom}}}
  {\sequence{\nat}}$ \\
  $= \arrlit
  {\sequence{\EraseA{\atval}}\,\dsum{\sequence{\idx}}{\fillc{\EraseC{\ctxt'}}{\EraseE{\expr}}}{}\,\sequence{\EraseA{\atom}}}
  {\sequence{\nat}}$ \\
  $= \fillc{\arrlit
    {\sequence{\EraseA{\atval}}\,\dsum{\sequence{\idx}}{\EraseC{\ctxt'}}{}\,\sequence{\EraseA{\atom}}}
    {\sequence{\nat}}}
  {\EraseE{\expr}}$ \\
  $= \fillc{\EraseC{\arrlit
    {\sequence{\atval}\,\dsum{\sequence{\idx}}{\ctxt'}{\type}\,\sequence{\atom}}
    {\sequence{\nat}}}}
  {\EraseE{\expr}}$ \\

\paragraph{Frame containing unevaluated cells:}
  \[\ctxt = \notedfrm[\type_r]{\sequence{\val}\,\ctxt'\,\sequence{\expr'}}{\sequence{\nat}}\]
  $\EraseE{\fillc{\notedfrm[\type_r]{\sequence{\val}\,\ctxt'\,\sequence{\expr}}{\sequence{\nat}}}{\expr}}$ \\
  $= \EraseE{\notedfrm[\type_r]{\sequence{\val}\,\fillc{\ctxt'}{\expr}\,\sequence{\expr'}}{\sequence{\nat}}}$ \\
  $= \frm{\sequence{\EraseE{\val}}\,\EraseE{\fillc{\ctxt'}{\expr}}\,\sequence{\EraseE{\expr'}}}{\EraseT{\type_r}}$ \\
  $= \frm{\sequence{\EraseE{\val}}\,\fillc{\EraseC{\ctxt'}}{\EraseE{\expr}}\,\sequence{\EraseE{\expr'}}}{\EraseT{\type_r}}$ \\
  $= \fillc{\frm{\sequence{\EraseE{\val}}\,\EraseC{\ctxt'}\,\sequence{\EraseE{\expr'}}}{\EraseT{\type_r}}}{\EraseE{\expr}}$ \\
  $= \fillc{\EraseC{\notedfrm[\type_r]
      {\sequence{\val}\,\EraseC{\ctxt'}\,\sequence{\expr'}}
      {\sequence{\nat}}}}
  {\EraseE{\expr}}$ \\

\paragraph{Application with unevaluated function:}
  \[\ctxt =
  \notedapp[\type_r]
  {\annotate[\typearray{\typefun{\sequence{\type_i}}{\type_o}}{\idx_f}]{\ctxt'}}
  {\sequence{\expr_a}}\]
  $\EraseE{\fillc
    {\notedapp[\type_r]
      {\annotate[\typearray{\typefun{\sequence{\type_i}}{\type_o}}{\idx_f}]{\ctxt'}}
      {\sequence{\expr_a}}}
    {\expr}}$\\
  $= \EraseE{\notedapp[\type_r]
    {\fillc{\annotate[\typearray{\typefun{\sequence{\type_i}}{\type_o}}{\idx_f}]{\ctxt'}}{\expr}}
    {\sequence{\expr_a}}}$ \\
  $= \app
  {\EraseE{\fillc{\ctxt'}{\expr}}}
  {\sequence{\ttsexp{\EraseE{\expr_a}}{\EraseT{\type_i}}}\;\EraseT{\type_r}}$ \\
  $= \app
  {\fillc{\EraseC{\ctxt'}}{\expr}}
  {\sequence{\ttsexp{\EraseE{\expr_a}}{\EraseT{\type_i}}}\;\EraseT{\type_r}}$ \\
  $= \fillc{\app
    {\EraseC{\ctxt'}}
    {\sequence{\ttsexp{\EraseE{\expr_a}}{\EraseT{\type_i}}}\;\EraseT{\type_r}}}
  {\expr}$ \\
  $= \fillc{\EraseC{\app
      {\annotate[\typearray{\typefun{\sequence{\type_i}}{\type_o}}{\idx_f}]{\ctxt'}}
      {\sequence{\expr_a}}}}
  {\expr}$ \\

\paragraph{Application with unevaluated argument:}
  \[\ctxt =
  \notedapp[\type_r]
  {\annotate
    [\typearray{\typefun{\sequence{\type_1}\,\type_2\,\sequence{\type_3}}{\type_o}}{\idx_f}]
    {\expr_f}}
  {\sequence{\val_1}\,\ctxt'\,\sequence{\expr_3}}\]
  $\EraseE{\fillc
    {\notedapp[\type_r]
      {\annotate
        [\typearray{\typefun{\sequence{\type_1}\,\type_2\,\sequence{\type_3}}{\type_o}}{\idx_f}]
        {\expr_f}}
      {\sequence{\val_1}\,\ctxt'\,\sequence{\expr_3}}}
    {\expr}}$ \\
  $= \EraseE{\notedapp[\type_r]
    {\annotate
      [\typearray{\typefun{\sequence{\type_1}\,\type_2\,\sequence{\type_3}}{\type_o}}{\idx_f}]
      {\expr_f}}
    {\sequence{\val_1}\,\fillc{\ctxt'}{\expr}\,\sequence{\expr_3}}}$ \\
  $= \app
  {\EraseE{\expr_f}}
  {\sequence{\ttsexp{\EraseE{\val_1}}{\EraseT{\type_1}}}
    \,\ttsexp{\EraseE{\fillc{\ctxt'}{\expr}}}{\EraseT{\type_2}}
    \,\sequence{\ttsexp{\EraseE{\expr_3}}{\EraseT{\type_3}}}
    \;\EraseT{\type_r}}$ \\
  $= \app
  {\EraseE{\expr_f}}
  {\sequence{\ttsexp{\EraseE{\val_1}}{\EraseT{\type_1}}}
    \,\ttsexp{\fillc{\EraseC{\ctxt'}}{\EraseE{\expr}}}{\EraseT{\type_2}}
    \,\sequence{\ttsexp{\EraseE{\expr_3}}{\EraseT{\type_3}}}
    \;\EraseT{\type_r}}$ \\
  $= \fillc
  {\app
    {\EraseE{\expr_f}}
    {\sequence{\ttsexp{\EraseE{\val_1}}{\EraseT{\type_1}}}
      \,\ttsexp{\EraseC{\ctxt'}}{\EraseT{\type_2}}
      \,\sequence{\ttsexp{\EraseE{\expr_3}}{\EraseT{\type_3}}}
      \;\EraseT{\type_r}}}
  {\EraseE{\expr}}$ \\
  $= \fillc
  {\EraseC{\notedapp[\type_r]
      {\annotate
        [\typearray{\typefun{\sequence{\type_1}\,\type_2\,\sequence{\type_3}}{\type_o}}{\idx_f}]
        {\expr_f}}
      {\sequence{\val_1}\,\ctxt'\,\sequence{\expr_3}}}}
  {\EraseE{\expr}}$ \\

\paragraph{Type application:}
  \[\ctxt =
  \notedtapp[\type_r]
  {\ctxt'}
  {\sequence{\type_a}}\]
  $\EraseE{\fillc{\notedtapp[\type_r]
      {\ctxt'}
      {\sequence{\type_a}}}
    {\expr}}$ \\
  $\EraseE{\notedtapp[\type_r]
    {\fillc{\ctxt'}{\expr}}
    {\sequence{\type_a}}}$ \\
  $= \iapp
  {\EraseE{\fillc{\ctxt'}{\expr}}}
  {\sequence{\EraseT{\type_a}}\;\EraseT{\type_r}}$ \\
  $= \iapp
  {\fillc{\EraseC{\ctxt'}}{\EraseE{\expr}}}
  {\sequence{\EraseT{\type_a}}\;\EraseT{\type_r}}$ \\
  $= \fillc{\iapp
    {\EraseC{\ctxt'}}
    {\sequence{\EraseT{\type_a}}\;\EraseT{\type_r}}}
  {\EraseE{\expr}}$ \\
  $= \fillc{\EraseC{\notedtapp[\type_r]
      {\ctxt'}
      {\sequence{\type_a}}}}
  {\EraseE{\expr}}$ \\

\paragraph{Index application:}
  \[\ctxt =
  \notediapp[\type_r]
  {\ctxt'}
  {\sequence{\idx_a}}\]
  $\EraseE{\fillc
    {\notediapp[\type_r]
      {\ctxt'}
      {\sequence{\idx_a}}}
    {\expr}}$ \\
  $= \EraseE{\notediapp[\type_r]
    {\fillc{\ctxt'}{\expr}}
    {\sequence{\idx_a}}}$ \\
  $= \iapp
  {\EraseE{\fillc{\ctxt'}{\expr}}}
  {\sequence{\idx_a}\;\EraseT{\type_r}}$ \\
  $= \iapp
  {\fillc{\EraseC{\ctxt'}}{\EraseE{\expr}}}
  {\sequence{\idx_a}\;\EraseT{\type_r}}$ \\
  $= \fillc{\iapp
    {\EraseC{\ctxt'}}
    {\sequence{\idx_a}\;\EraseT{\type_r}}}
  {\EraseE{\expr}}$ \\
  $= \fillc{\EraseC{\notediapp[\type_r]
      {\ctxt'}
      {\sequence{\idx_a}}}}
  {\EraseE{\expr}}$ \\

\paragraph{Unboxing:}
  \[\ctxt =
  \dproj
  {\sequence{\var_i}}
  {\var_e}
  {\ctxt'}
  {\expr_b^{\type_b}}\]
  $\EraseE{\fillc
    {\dproj
      {\sequence{\var_i}}
      {\var_e}
      {\ctxt'}
      {\expr_b^{\type_b}}}
    {\expr}}$ \\
  $= \EraseE{\dproj
    {\sequence{\var_i}}
    {\var_e}
    {\fillc{\ctxt'}{\expr}}
    {\expr_b^{\type_b}}}$ \\
  $= \dproj
  {\sequence{\var_i}}
  {\var_e}
  {\EraseE{\fillc{\ctxt'}{\expr}}}
  {\EraseE{\expr_b^{\type_b}} \; \EraseT{\type_b}}$ \\
  $= \dproj
  {\sequence{\var_i}}
  {\var_e}
  {\fillc{\EraseC{\ctxt'}}{\EraseE{\expr}}}
  {\EraseE{\expr_b^{\type_b}} \; \EraseT{\type_b}}$ \\
  $= \fillc{\dproj
    {\sequence{\var_i}}
    {\var_e}
    {\EraseC{\ctxt'}}
    {\EraseE{\expr_b^{\type_b}} \; \EraseT{\type_b}}}
  {\EraseE{\expr}}$ \\
  $= \fillc{\EraseC{\dproj
      {\sequence{\var_i}}
      {\var_e}
      {\ctxt'}
      {\expr_b}}}
  {\EraseE{\expr}}$ \\
\end{sproof}

We will also rely on a series of lemmas showing that substitution commutes with erasure.
% Erase a generic expr then sub it into something
\newcommand{\eexprsubst}[1]{\subst{#1}{\var}{\EraseE{\expr_\var}}}
\newcommand{\etrmsubstE}[1]{\eexprsubst{\EraseTrm{#1}}}
\newcommand{\exprsubstE}[1]{\eexprsubst{\EraseE{#1}}}
\newcommand{\eatomsubstE}[1]{\eexprsubst{\EraseA{#1}}}
\begin{lemma}%[Term substitution commutes with erasure]
  \label{EESubstErase}
  $\EraseTrm{\eexprsubst{\term}} = \etrmsubstE{\term}$
\end{lemma}
\begin{sproof}[This is straightforward induction on $\term$.]
We use induction on $\term$.
We skip the cases where $\var$ does not appear free in $\term$,
as substitution would not change $\term$.
\paragraph{Term abstraction:}
\[\term = \lam{\sequence{\notevar{\var_i}{\type}}}{\expr}\text{, where } \var \not\in \sequence{\var_i}\]
$\EraseTrm{\eexprsubst{\lam{\sequence{\notevar{\var_i}{\type}}}{\expr}}}$ \\
$= \EraseA{\eexprsubst{\lam{\sequence{\notevar{\var_i}{\type}}}{\expr}}}$ \\
$= \EraseA{\lam{\sequence{\notevar{\var_i}{\type}}}{\eexprsubst{\expr}}}$ \\
$= \lam{\sequence{\var_i}}{\EraseE{\eexprsubst{\expr}}}$ \\
$= \lam{\sequence{\var_i}}{\exprsubstE{\expr}}$ \\
$= \eexprsubst{\lam{\sequence{\var_i}}{\EraseE{\expr}}}$ \\
$= \eatomsubstE{\lam{\sequence{\notevar{\var_i}{\type}}}{\expr}}$ \\
$= \etrmsubstE{\lam{\sequence{\notevar{\var_i}{\type}}}{\expr}}$ \\

\paragraph{Type abstraction:}
\[\term = \tlam{\sequence{\notevar{\var_i}{\kind}}}{\val}\]
$\EraseTrm{\eexprsubst{\tlam{\sequence{\notevar{\var_i}{\kind}}}{\val}}}$ \\
$= \EraseA{\eexprsubst{\tlam{\sequence{\notevar{\var_i}{\kind}}}{\val}}}$ \\
$= \EraseA{\tlam{\sequence{\notevar{\var_i}{\kind}}}{\eexprsubst{\val}}}$ \\
$= \ilam{\sequence{\var_i}}{\EraseE{\eexprsubst{\val}}}$ \\
$= \ilam{\sequence{\var_i}}{\exprsubstE{\val}}$ \\
$= \eexprsubst{\ilam{\sequence{\var_i}}{\EraseE{\val}}}$ \\
$= \eatomsubstE{\tlam{\sequence{\notevar{\var_i}{\kind}}}{\val}}$ \\
$= \etrmsubstE{\tlam{\sequence{\notevar{\var_i}{\kind}}}{\val}}$ \\

\paragraph{Index abstraction:}
\[\term = \ilam{\sequence{\notevar{\var_i}{\sort}}}{\val}\]
$\EraseTrm{\eexprsubst{\ilam{\sequence{\notevar{\var_i}{\sort}}}{\val}}}$ \\
$= \EraseA{\eexprsubst{\ilam{\sequence{\notevar{\var_i}{\sort}}}{\val}}}$ \\
$= \EraseA{\ilam{\sequence{\notevar{\var_i}{\sort}}}{\eexprsubst{\val}}}$ \\
$= \ilam{\sequence{\var_i}}{\EraseE{\eexprsubst{\val}}}$ \\
$= \ilam{\sequence{\var_i}}{\exprsubstE{\val}}$ \\
$= \eexprsubst{\ilam{\sequence{\var_i}}{\EraseE{\val}}}$ \\
$= \eatomsubstE{\ilam{\sequence{\notevar{\var_i}{\sort}}}{\val}}$ \\
$= \etrmsubstE{\ilam{\sequence{\notevar{\var_i}{\sort}}}{\val}}$ \\

\paragraph{Box:}
\[\term = \dsum{\sequence{\idx}}{\expr_s}{\type}\]
$\EraseTrm{\eexprsubst{\dsum{\sequence{\idx}}{\expr_s}{\type}}}$ \\
$= \EraseA{\eexprsubst{\dsum{\sequence{\idx}}{\expr_s}{\type}}}$ \\
$= \EraseA{\dsum{\sequence{\idx}}{\eexprsubst{\expr_s}}{\type}}$ \\
$= \dsum{\sequence{\idx}}{\EraseE{\eexprsubst{\expr_s}}}{}$ \\
$= \dsum{\sequence{\idx}}{\exprsubstE{\expr_s}}{}$ \\
$= \eexprsubst{\dsum{\sequence{\idx}}{\EraseE{\expr_s}}{}}$ \\
$= \eatomsubstE{\dsum{\sequence{\idx}}{\expr_s}{\type}}$ \\
$= \etrmsubstE{\dsum{\sequence{\idx}}{\expr_s}{\type}}$ \\

\paragraph{Variable:}
\[\term = \var\]
$\EraseTrm{\eexprsubst{\var}}$ \\
$= \EraseE{\eexprsubst{\var}}$ \\
$= \EraseE{\expr_\var}$ \\
$= \eexprsubst{\var}$ \\
$= \exprsubstE{\var}$ \\
$= \etrmsubstE{\var}$ \\

\paragraph{Array:}
\[\term = \arrlit{\sequence{\atom}}{\sequence{\nat}}\]
$\EraseTrm{\eexprsubst{\arrlit{\sequence{\atom}}{\sequence{\nat}}}}$ \\
$= \EraseE{\eexprsubst{\arrlit{\sequence{\atom}}{\sequence{\nat}}}}$ \\
$= \EraseE{\arrlit{\sequence{\eexprsubst{\atom}}}{\sequence{\nat}}}$ \\
$= \arrlit{\sequence{\EraseA{\eexprsubst{\atom}}}}{\sequence{\nat}}$ \\
$= \arrlit{\sequence{\eatomsubstE{\atom}}}{\sequence{\nat}}$ \\
$= \eexprsubst{\arrlit{\sequence{\EraseA{\atom}}}{\sequence{\nat}}}$ \\
$= \exprsubstE{\arrlit{\sequence{\atom}}{\sequence{\nat}}}$ \\
$= \etrmsubstE{\arrlit{\sequence{\atom}}{\sequence{\nat}}}$ \\

\paragraph{Frame:}
\[\term = \notedfrm[\type_r]{\sequence{\expr_c}}{\sequence{\nat}}\]
$\EraseTrm{\eexprsubst{\notedfrm[\type_r]{\sequence{\expr_c}}{\sequence{\nat}}}}$ \\
$= \EraseE{\eexprsubst{\notedfrm[\type_r]{\sequence{\expr_c}}{\sequence{\nat}}}}$ \\
$= \EraseE{\notedfrm[\type_r]{\sequence{\eexprsubst{\expr_c}}}{\sequence{\nat}}}$ \\
$= \frm{\sequence{\EraseE{\eexprsubst{\expr_c}}}}{\EraseT{\type_r}}$ \\
$= \frm{\sequence{\exprsubstE{\expr_c}}}{\EraseT{\type_r}}$ \\
$= \eexprsubst{\frm{\sequence{\EraseE{\expr_c}}}{\EraseT{\type_r}}}$ \\
$= \exprsubstE{\notedfrm[\type_r]{\sequence{\expr_c}}{\sequence{\nat}}}$ \\
$= \etrmsubstE{\notedfrm[\type_r]{\sequence{\expr_c}}{\sequence{\nat}}}$ \\

\paragraph{Application:}
\[\term =\notedapp[\type_r]
{\annotate
  [\typearray{\typefun{\sequence{\type_i}}{\type_o}}{\idx_f}]
  {\expr_f}}
{\sequence{\expr_a}}\]
$\EraseTrm{\eexprsubst{\notedapp[\type_r]{\annotate
      [\typearray{\typefun{\sequence{\type_i}}{\type_o}}{\idx_f}]
      {\expr_f}}
    {\sequence{\expr_a}}}}$ \\
$= \EraseE{\eexprsubst{\notedapp[\type_r]{\annotate
      [\typearray{\typefun{\sequence{\type_i}}{\type_o}}{\idx_f}]
      {\expr_f}}
    {\sequence{\expr_a}}}}$ \\
$= \EraseE{\notedapp[\type_r]{\eexprsubst{\annotate
      [\typearray{\typefun{\sequence{\type_i}}{\type_o}}{\idx_f}]
      {\expr_f}}}
  {\sequence{\eexprsubst{\expr_a}}}}$ \\
$= \app
{\EraseE{\eexprsubst{\expr_f}}}
{\sequence{\ttsexp{\EraseE{\eexprsubst{\expr_a}}}{\EraseT{\type_i}}}\;\EraseT{\type_r}}$ \\
$= \app
{\exprsubstE{\expr_f}}
{\sequence{\ttsexp{\exprsubstE{\expr_a}}{\EraseT{\type_i}}}\;\EraseT{\type_r}}$ \\
$= \eexprsubst{\app
  {\EraseE{\expr_f}}
  {\sequence{\ttsexp{\EraseE{\expr_a}}{\EraseT{\type_i}}}\;\EraseT{\type_r}}}$ \\
$= \exprsubstE{\notedapp[\type_r]
  {\annotate
    [\typearray{\typefun{\sequence{\type_i}}{\type_o}}{\idx_f}]
    {\expr_f}}
  {\sequence{\expr_a}}}$ \\
$= \etrmsubstE{\notedapp[\type_r]
  {\annotate
    [\typearray{\typefun{\sequence{\type_i}}{\type_o}}{\idx_f}]
    {\expr_f}}
  {\sequence{\expr_a}}}$ \\

\paragraph{Type application:}
\[\term = \notedtapp[\type_r]{\expr_f}{\sequence{\type_a}}\]
$\EraseTrm{\eexprsubst{\notedtapp[\type_r]{\expr_f}{\sequence{\type_a}}}}$ \\
$= \EraseE{\eexprsubst{\notedtapp[\type_r]{\expr_f}{\sequence{\type_a}}}}$ \\
$= \EraseE{\notedtapp[\type_r]{\eexprsubst{\expr_f}}{\sequence{\type_a}}}$ \\
$= \iapp{\EraseE{\eexprsubst{\expr_f}}}{\sequence{\EraseT{\type_a}}\;\EraseT{\type_r}}$ \\
$= \iapp{\exprsubstE{\expr_f}}{\sequence{\EraseT{\type_a}}\;\EraseT{\type_r}}$ \\
$= \eexprsubst{\iapp{\EraseE{\expr_f}}{\sequence{\EraseT{\type_a}}\;\EraseT{\type_r}}}$ \\
$= \exprsubstE{\notedtapp[\type_r]{\expr_f}{\sequence{\type_a}}}$ \\
$= \etrmsubstE{\notedtapp[\type_r]{\expr_f}{\sequence{\type_a}}}$ \\

\paragraph{Index application:}
\[\term = \notediapp[\type_r]{\expr_f}{\sequence{\idx_a}}\]
$\EraseTrm{\eexprsubst{\notediapp[\type_r]{\expr_f}{\sequence{\idx_a}}}}$ \\
$= \EraseE{\eexprsubst{\notediapp[\type_r]{\expr_f}{\sequence{\idx_a}}}}$ \\
$= \EraseE{\notediapp[\type_r]{\eexprsubst{\expr_f}}{\sequence{\idx_a}}}$ \\
$= \iapp{\EraseE{\eexprsubst{\expr_f}}}{\sequence{\idx_a}\;\EraseT{\type_r}}$ \\
$= \iapp{\exprsubstE{\expr_f}}{\sequence{\idx_a}\;\EraseT{\type_r}}$ \\
$= \eexprsubst{\iapp{\EraseE{\expr_f}}{\sequence{\idx_a}\;\EraseT{\type_r}}}$ \\
$= \exprsubstE{\notediapp[\type_r]{\expr_f}{\sequence{\idx_a}}}$ \\
$= \etrmsubstE{\notediapp[\type_r]{\expr_f}{\sequence{\idx_a}}}$ \\

\paragraph{Unboxing:}
\[\term = \dproj{\sequence{\var_i}}{\var_e}{\expr_s}{\expr_b^{\type_b}}\]
$\EraseTrm{\eexprsubst{\dproj{\sequence{\var_i}}{\var_e}{\expr_s}{\expr_b^{\type_b}}}}$ \\
$= \EraseE{\eexprsubst{\dproj{\sequence{\var_i}}{\var_e}{\expr_s}{\expr_b^{\type_b}}}}$ \\
$= \EraseE{\dproj{\sequence{\var_i}}{\var_e}{\eexprsubst{\expr_s}}{\eexprsubst{\expr_b^{\type_b}}}}$ \\
$= \dproj{\sequence{\var_i}}{\var_e}{\EraseE{\eexprsubst{\expr_s}}}
{\EraseE{\eexprsubst{\expr_b^{\type_b}}}\;\EraseT{\type_b}}$ \\
% Note: the above line relies on an earlier lemma that substitution in e_b does not change its type
$= \dproj{\sequence{\var_i}}{\var_e}{\exprsubstE{\expr_s}}{\exprsubstE{\expr_b^{\type_b}}\;\EraseT{\type_b}}$ \\
$= \eexprsubst{\dproj{\sequence{\var_i}}{\var_e}{\expr_s}{\expr_b\;\EraseT{\type_b}}}$ \\
$= \exprsubstE{\dproj{\sequence{\var_i}}{\var_e}{\expr_s}{\expr_b}}$ \\
$= \etrmsubstE{\dproj{\sequence{\var_i}}{\var_e}{\expr_s}{\expr_b}}$ \\

\end{sproof}

\newcommand{\etypesubst}[1]{\subst{#1}{\var}{\EraseT{\type_\var}}}
\begin{lemma}%[Type-into-type substitution commutes with erasure]
  \label{TTSubstErase}
  $\EraseT{\typesubst{\type}} = \etypesubst{\EraseT{\type}}$
\end{lemma}
\begin{sproof}[This is straightforward induction on $\type$.]
  We use induction on $\type$.
  We elide the cases where $\var$ does not appear free in $\type$.
  \paragraph{Variable:}
  \[\type = \var\]
  $\EraseT{\typesubst{\var}}$ \\
  $= \EraseT{\type_\var}$ \\
  $= \etypesubst{\var}$ \\
  $= \etypesubst{\EraseT{\var}}$ \\

  \paragraph{Function:}
  \[\type = \typefun{\sequence{\type_i}}{\type_o}\]
  $\EraseT{\typesubst{\typefun{\sequence{\type_i}}{\type_o}}}$ \\
  $= \EraseT{\typefun{\sequence{\typesubst{\type_i}}}{\typesubst{\type_o}}}$ \\
  $= \idxshape{}$ \\
  % $= \etypesubst{\idxshape{}}$ \\
  $= \etypesubst{\EraseT{\typefun{\sequence{\type_i}}{\type_o}}}$ \\

  \paragraph{Universal:}
  \[\type = \typeuniv{\sequence{\notevar{\var_u}{\kind}}}{\type_u}\text{, where }\var \not\in \sequence{\var_u}\]
  $\EraseT{\typesubst{\typeuniv{\sequence{\notevar{\var_u}{\kind}}}{\type_u}}}$ \\
  $= \EraseT{\typeuniv{\sequence{\notevar{\var_u}{\kind}}}{\typesubst{\type_u}}}$ \\
  $= \idxshape{}$ \\
  $= \etypesubst{\EraseT{\typeuniv{\sequence{\notevar{\var_u}{\kind}}}{\type_u}}}$

  \paragraph{Dependent product:}
  \[\type = \typedprod{\sequence{\notevar{\var_p}{\sort}}}{\type_p}\]
  $\EraseT{\typesubst{\typedprod{\sequence{\notevar{\var_p}{\sort}}}{\type_p}}}$ \\
  $= \EraseT{\typedprod{\sequence{\notevar{\var_p}{\sort}}}{\typesubst{\type_p}}}$ \\
  $= \idxshape{}$ \\
  $= \etypesubst{\EraseT{\typedprod{\sequence{\notevar{\var_p}{\sort}}}{\type_p}}}$ \\

  \paragraph{Dependent sum:}
  \[\type = \typedsum{\sequence{\notevar{\var_p}{\sort}}}{\type_p}\]
  $\EraseT{\typesubst{\typedsum{\sequence{\notevar{\var_p}{\sort}}}{\type_p}}}$ \\
  $= \EraseT{\typedsum{\sequence{\notevar{\var_p}{\sort}}}{\typesubst{\type_p}}}$ \\
  $= \idxshape{}$ \\
  $= \etypesubst{\EraseT{\typedsum{\sequence{\notevar{\var_p}{\sort}}}{\type_p}}}$ \\

  \paragraph{Array:}
  \[\type = \typearray{\type_a}{\idx}\]
  $\EraseT{\typesubst{\typearray{\type_a}{\idx}}}$ \\
  $= \EraseT{\typearray{\typesubst{\type_a}}{\idx_a}}$ \\
  $= \idx_a$\qquad\text{(note: $\var$ is a type variable and cannot appear in $\idx_a$)} \\
  $= \etypesubst{\EraseT{\typearray{\type_a}{\idx}}}$

\end{sproof}

%\newcommand{\tesubstE}[1]{\subst}
\begin{lemma}%[Type-into-term substitution commutes with erasure]
  \label{TESubstErase}
  $\EraseTrm{\typesubst{\term}} = \etypesubst{\EraseTrm{\term}}$
\end{lemma}
\begin{sproof}[This is straightforward induction on $\term$.]
  We use induction on $\term$.
  We elide the cases where $\var$ does not appear free in $\term$.

  \paragraph{Term abstraction:}
  \[\term = \lam{\sequence{\notevar{\var_i}{\type}}}{\expr}\]
  $\EraseTrm{\typesubst{\lam{\sequence{\notevar{\var_i}{\type}}}{\expr}}}$ \\
  $= \EraseA{\typesubst{\lam{\sequence{\notevar{\var_i}{\type}}}{\expr}}}$ \\
  $= \EraseA{\lam{\sequence{\notevar{\var_i}{\type}}}{\typesubst{\expr}}}$ \\
  $= \lam{\sequence{\var_i}}{\EraseE{\etypesubst{\expr}}}$ \\
  $= \lam{\sequence{\var_i}}{\etypesubst{\EraseE{\expr}}}$ \\
  $= \etypesubst{\lam{\sequence{\var_i}}{\EraseE{\expr}}}$ \\
  $= \etypesubst{\EraseA{\lam{\sequence{\notevar{\var_i}{\type}}}{\expr}}}$ \\
  $= \etypesubst{\EraseTrm{\lam{\sequence{\notevar{\var_i}{\type}}}{\expr}}}$ \\

  \paragraph{Type abstraction:}
  \[\term = \tlam{\sequence{\notevar{\var_i}{\kind}}}{\val}\text{, where } \var \not\in \sequence{\var_i}\]
  $\EraseTrm{\typesubst{\tlam{\sequence{\notevar{\var_i}{\kind}}}{\val}}}$ \\
  $= \EraseA{\typesubst{\tlam{\sequence{\notevar{\var_i}{\kind}}}{\val}}}$ \\
  $= \EraseA{\tlam{\sequence{\notevar{\var_i}{\kind}}}{\typesubst{\val}}}$ \\
  $= \ilam{\sequence{\var_i}}{\EraseE{\typesubst{\val}}}$ \\
  $= \ilam{\sequence{\var_i}}{\etypesubst{\EraseE{\val}}}$ \\
  $= \etypesubst{\ilam{\sequence{\var_i}}{\EraseE{\val}}}$ \\
  $= \etypesubst{\EraseA{\tlam{\sequence{\notevar{\var_i}{\kind}}}{\val}}}$ \\
  $= \etypesubst{\EraseTrm{\tlam{\sequence{\notevar{\var_i}{\kind}}}{\val}}}$ \\

  \paragraph{Index abstraction:}
  \[\term = \ilam{\sequence{\notevar{\var_i}{\sort}}}{\val}\]
  $\EraseTrm{\typesubst{\ilam{\sequence{\notevar{\var_i}{\sort}}}{\val}}}$ \\
  $= \EraseA{\typesubst{\ilam{\sequence{\notevar{\var_i}{\sort}}}{\val}}}$ \\
  $= \EraseA{\ilam{\sequence{\notevar{\var_i}{\sort}}}{\typesubst{\val}}}$ \\
  $= \ilam{\sequence{\var_i}}{\EraseE{\typesubst{\val}}}$ \\
  $= \ilam{\sequence{\var_i}}{\etypesubst{\EraseE{\val}}}$ \\
  $= \etypesubst{\ilam{\sequence{\var_i}}{\EraseE{\val}}}$ \\
  $= \etypesubst{\EraseA{\ilam{\sequence{\notevar{\var_i}{\kind}}}{\val}}}$ \\
  $= \etypesubst{\EraseTrm{\ilam{\sequence{\notevar{\var_i}{\kind}}}{\val}}}$ \\

  \paragraph{Box:}
  \[\term = \dsum{\sequence{\idx}}{\expr_s}{\type}\]
  $\EraseTrm{\typesubst{\dsum{\sequence{\idx}}{\expr_s}{\type}}}$ \\
  $= \EraseA{\typesubst{\dsum{\sequence{\idx}}{\expr_s}{\type}}}$ \\
  $= \EraseA{{\dsum{\sequence{\idx}}{\typesubst{\expr_s}}{\typesubst{\type}}}}$ \\
  $= \dsum{\sequence{\idx}}{\EraseE{\typesubst{\expr_s}}}{}$ \\
  $= \dsum{\sequence{\idx}}{\etypesubst{\EraseE{\expr_s}}}{}$ \\
  $= \etypesubst{\dsum{\sequence{\idx}}{\EraseE{\expr_s}}{}}$ \\
  $= \etypesubst{\EraseA{\dsum{\sequence{\idx}}{\expr_s}{\type}}}$ \\
  $= \etypesubst{\EraseTrm{\dsum{\sequence{\idx}}{\expr_s}{\type}}}$ \\

  \paragraph{Array:}
  \[\term = \arrlit{\sequence{\atom}}{\sequence{\nat}}\]
  $\EraseTrm{\typesubst{\arrlit{\sequence{\atom}}{\sequence{\nat}}}}$ \\
  $= \EraseE{\typesubst{\arrlit{\sequence{\atom}}{\sequence{\nat}}}}$ \\
  $= \EraseE{\arrlit{\sequence{\typesubst{\atom}}}{\sequence{\nat}}}$ \\
  $= \arrlit{\sequence{\EraseA{\typesubst{\atom}}}}{\sequence{\nat}}$ \\
  $= \arrlit{\sequence{\etypesubst{\EraseA{\atom}}}}{\sequence{\nat}}$ \\
  $= \etypesubst{\arrlit{\sequence{\EraseA{\atom}}}{\sequence{\nat}}}$ \\
  $= \etypesubst{\EraseE{\arrlit{\sequence{\atom}}{\sequence{\nat}}}}$ \\
  $= \etypesubst{\EraseTrm{\arrlit{\sequence{\atom}}{\sequence{\nat}}}}$ \\

  \paragraph{Frame:}
  \[\term = \notedfrm[\type_r]{\sequence{\expr_c}}{\sequence{\nat}}\]
  $\EraseTrm{\typesubst{\notedfrm[\type_r]{\sequence{\expr_c}}{\sequence{\nat}}}}$ \\
  $= \EraseE{\typesubst{\notedfrm[\type_r]{\sequence{\expr_c}}{\sequence{\nat}}}}$ \\
  $= \EraseE{\notedfrm[\typesubst{\type_r}]{\sequence{\typesubst{\expr_c}}}{\sequence{\nat}}}$ \\
  $= \frm{\sequence{\EraseE{\typesubst{\expr_c}}}}{\EraseT{\typesubst{\type_r}}}$ \\
  $= \frm{\sequence{\etypesubst{\EraseE{\expr_c}}}}{\EraseT{\typesubst{\type_r}}}$%
  , by the induction hypothesis \\
  $= \frm{\sequence{\etypesubst{\EraseE{\expr_c}}}}{\etypesubst{\EraseT{\type_r}}}$%
  , by Lemma \ref{TTSubstErase} \\
  $= \etypesubst{\frm{\sequence{\EraseE{\expr_c}}}{\EraseT{\type_r}}}$ \\
  $= \etypesubst{\EraseE{\notedfrm[\type_r]{\sequence{\expr_c}}{\sequence{\nat}}}}$ \\
  $= \etypesubst{\EraseTrm{\notedfrm[\type_r]{\sequence{\expr_c}}{\sequence{\nat}}}}$ \\

  \paragraph{Application:}
  \[\term =
  \notedapp[\type_r]{\annotate
    [\typearray{\typefun{\sequence{\type_i}}{\type_o}}{\idx_f}]
    {\expr_f}}{\sequence{\expr_a}}\]
  $\EraseTrm{\typesubst{\notedapp[\type_r]{\annotate
        [\typearray{\typefun{\sequence{\type_i}}{\type_o}}{\idx_f}]
        {\expr_f}}{\sequence{\expr_a}}}}$ \\
  $= \EraseE{\typesubst{\notedapp[\type_r]{\annotate
        [\typearray{\typefun{\sequence{\type_i}}{\type_o}}{\idx_f}]
        {\expr_f}}{\sequence{\expr_a}}}}$ \\
  $= \EraseE{\notedapp[\typesubst{\type_r}]{\annotate
      [\typesubst{\typearray{\typefun{\sequence{\type_i}}{\type_o}}{\idx_f}}]
      {\typesubst{\expr_f}}}{\sequence{\typesubst{\expr_a}}}}$ \\
  $= \EraseE{\notedapp[\typesubst{\type_r}]{\annotate
      [\typearray{\typefun{\sequence{\typesubst{\type_i}}}{\typesubst{\type_o}}}{\idx_f}]
      {\typesubst{\expr_f}}}{\sequence{\typesubst{\expr_a}}}}$ \\
  $= \app{\EraseE{\typesubst{\expr_f}}}{\sequence
    {\ttsexp{\EraseE{\typesubst{\expr_a}}}{\EraseT{\typesubst{\type_i}}}}
    \;\EraseT{\typesubst{\type_r}}}$ \\
  $= \app{\etypesubst{\EraseE{\expr_f}}}{\sequence
    {\ttsexp{\etypesubst{\EraseE{\expr_a}}}{\EraseT{\typesubst{\type_i}}}}
    \;\EraseT{\typesubst{\type_r}}}$, by the induction hypothesis \\
  $= \app{\etypesubst{\EraseE{\expr_f}}}{\sequence
    {\ttsexp{\etypesubst{\EraseE{\expr_a}}}{\etypesubst{\EraseT{\type_i}}}}
    \;\etypesubst{\EraseE{\type_r}}}$, by Lemma \ref{TTSubstErase} \\
  $= \etypesubst{\app{\EraseE{\expr_f}}{\sequence
    {\ttsexp{\EraseE{\expr_a}}{\EraseT{\type_i}}}
    \;\EraseE{\type_r}}}$ \\
  $= \etypesubst{\EraseE{\notedapp[\type_r]{\annotate
        [\typearray{\typefun{\sequence{\type_i}}{\type_o}}{\idx_f}]
        {\expr_f}}{\sequence{\expr_a}}}}$ \\
  $= \etypesubst{\EraseTrm{\notedapp[\type_r]{\annotate
        [\typearray{\typefun{\sequence{\type_i}}{\type_o}}{\idx_f}]
        {\expr_f}}{\sequence{\expr_a}}}}$ \\

  \paragraph{Type application:}
  \[\term = \notedtapp[\type_r]{\expr_f}{\sequence{\type_a}}\]
  $\EraseTrm{\typesubst{\notedtapp[\type_r]{\expr_f}{\sequence{\type_a}}}}$ \\
  $= \EraseE{\typesubst{\notedtapp[\type_r]{\expr_f}{\sequence{\type_a}}}}$ \\
  $= \EraseE{\notedtapp[\typesubst{\type_r}]{\typesubst{\expr_f}}{\sequence{\typesubst{\type_a}}}}$ \\
  $= \iapp{\EraseE{\typesubst{\expr_f}}}{\sequence
    {\EraseT{\typesubst{\type_a}}}\;\EraseT{\typesubst{\type_r}}}$ \\
  $= \iapp{\etypesubst{\EraseE{\expr_f}}}{\sequence
    {\EraseT{\typesubst{\type_a}}}\;\EraseT{\typesubst{\type_r}}}$%
  , by the induction hypothesis \\
  $= \iapp{\etypesubst{\EraseE{\expr_f}}}{\sequence
    {\etypesubst{\EraseT{\type_a}}}\;\etypesubst{\EraseT{\type_r}}}$%
  , by Lemma \ref{TTSubstErase} \\
  $= \etypesubst{\iapp{\EraseE{\expr_f}}{\sequence
      {\EraseT{\type_a}}\;\EraseT{\type_r}}}$ \\
  $= \etypesubst{\EraseE{\notedtapp[\type_r]{\expr_f}{\sequence{\type_a}}}}$ \\
  $= \etypesubst{\EraseTrm{\notedtapp[\type_r]{\expr_f}{\sequence{\type_a}}}}$ \\

  \paragraph{Index application:}
  \[\term = \notediapp[\type_r]{\expr_f}{\sequence{\idx_a}}\]
  $\EraseTrm{\typesubst{\notediapp[\type_r]{\expr_f}{\sequence{\idx_a}}}}$ \\
  $= \EraseE{\typesubst{\notediapp[\type_r]{\expr_f}{\sequence{\idx_a}}}}$ \\
  $= \EraseE{\notediapp[\typesubst{\type_r}]{\typesubst{\expr_f}}{\sequence{\idx_a}}}$ \\
  $= \iapp{\EraseE{\typesubst{\expr_f}}}{\sequence{\idx_a}\;\EraseT{\typesubst{\type_r}}}$ \\
  $= \iapp{\etypesubst{\EraseE{\expr_f}}}{\sequence{\idx_a}\;\EraseT{\typesubst{\type_r}}}$%
  , by the induction hypothesis\\
  $= \iapp{\etypesubst{\EraseE{\expr_f}}}{\sequence{\idx_a}\;\etypesubst{\EraseT{\type_r}}}$%
  , by Lemma \ref{TTSubstErase}\\
  $= \etypesubst{\iapp{\EraseE{\expr_f}}{\sequence{\idx_a}\;\EraseT{\type_r}}}$ \\
  $= \etypesubst{\EraseE{\notediapp[\type_r]{\expr_f}{\sequence{\idx_a}}}}$ \\
  $= \etypesubst{\EraseTrm{\notediapp[\type_r]{\expr_f}{\sequence{\idx_a}}}}$ \\

  \paragraph{Unboxing:}
  \[\term = \dproj{\sequence{\var_i}}{\var_e}{\expr_s}{\expr_b^{\type_b}}\]
  $\EraseTrm{\typesubst{\dproj{\sequence{\var_i}}{\var_e}{\expr_s}{\expr_b^{\type_b}}}}$ \\
  $= \EraseE{\typesubst{\dproj{\sequence{\var_i}}{\var_e}{\expr_s}{\expr_b^{\type_b}}}}$ \\
  $= \EraseE{\dproj{\sequence{\var_i}}{\var_e}{\typesubst{\expr_s}}
    {\typesubst{\expr_b}^{\typesubst{\type_b}}}}$ \\
  $= \dproj{\sequence{\var_i}}{\var_e}{\EraseE{\typesubst{\expr_s}}}
  {\EraseE{\typesubst{\expr_b}} \; \EraseT{\typesubst{\type_b}}}$ \\
  $= \dproj{\sequence{\var_i}}{\var_e}{\EraseE{\typesubst{\expr_s}}}
  {\EraseE{\typesubst{\expr_b}} \; \etypesubst{\EraseT{\type_b}}}$ \\
  $= \dproj{\sequence{\var_i}}{\var_e}{\etypesubst{\EraseE{\expr_s}}}
  {\etypesubst{\EraseE{\expr_b}} \; \etypesubst{\EraseT{\type_b}}}$ \\
  $= \etypesubst{\dproj{\sequence{\var_i}}{\var_e}{\EraseE{\expr_s}}{\EraseE{\expr_b} \; \EraseT{\type_b}}}$ \\
  $= \etypesubst{\EraseE{\dproj{\sequence{\var_i}}{\var_e}{\expr_s}{\expr_b}}}$ \\
  $= \etypesubst{\EraseTrm{\dproj{\sequence{\var_i}}{\var_e}{\expr_s}{\expr_b}}}$ \\

\end{sproof}

\begin{lemma}%[Index-into-type substitution commutes with erasure]
  \label{ITSubstErase}
  $\EraseT{\idxsubst{\type}} = \idxsubst{\EraseT{\type}}$
\end{lemma}
\begin{sproof}[This is straightforward induction on $\type$.]
  We use induction on $\type$.
  We elide the cases where $\var$ does not appear free in $\type$.
  \paragraph{Function:}
  \[\type = \typefun{\sequence{\type_i}}{\type_o}\]
  $\EraseT{\idxsubst{\typefun{\sequence{\type_i}}{\type_o}}}$ \\
  $= \EraseT{\typefun{\sequence{\idxsubst{\type_i}}}{\idxsubst{\type_o}}}$ \\
  $= \idxshape{}$ \\
  $= \idxsubst{\idxshape{}}$ \\
  $= \idxsubst{\EraseT{\typefun{\sequence{\type_i}}{\type_o}}}$ \\

  \paragraph{Universal:}
  \[\type = \typeuniv{\sequence{\notevar{\var_u}{\kind}}}{\type_u}\]
  $\EraseT{\idxsubst{\typeuniv{\sequence{\notevar{\var_u}{\kind}}}{\type_u}}}$ \\
  $= \EraseT{\typeuniv{\sequence{\notevar{\var_u}{\kind}}}{\idxsubst{\type_u}}}$ \\
  $= \idxshape{}$ \\
  $= \idxsubst{\idxshape{}}$ \\
  $= \idxsubst{\EraseT{\typeuniv{\sequence{\notevar{\var_u}{\kind}}}{\type_u}}}$ \\

  \paragraph{Dependent product:}
  \[\type = \typedprod{\sequence{\notevar{\var_p}{\sort}}}{\type_p}\text{, where }\var \not\in \sequence{\var_p}\]
  $\EraseT{\idxsubst{\typedprod{\sequence{\notevar{\var_p}{\sort}}}{\type_p}}}$ \\
  $= \EraseT{\typedprod{\sequence{\notevar{\var_p}{\sort}}}{\idxsubst{\type_p}}}$ \\
  $= \idxshape{}$ \\
  $= \idxsubst{\idxshape{}}$ \\
  $= \idxsubst{\EraseT{\typedprod{\sequence{\notevar{\var_p}{\sort}}}{\type_p}}}$ \\

  \paragraph{Dependent sum:}
  \[\type = \typedsum{\sequence{\notevar{\var_p}{\sort}}}{\type_p}\text{, where }\var \not\in \sequence{\var_p}\]
  $\EraseT{\idxsubst{\typedsum{\sequence{\notevar{\var_p}{\sort}}}{\type_p}}}$ \\
  $= \EraseT{\typedsum{\sequence{\notevar{\var_p}{\sort}}}{\idxsubst{\type_p}}}$ \\
  $= \idxshape{}$ \\
  $= \idxsubst{\idxshape{}}$ \\
  $= \idxsubst{\EraseT{\typedsum{\sequence{\notevar{\var_p}{\sort}}}{\type_p}}}$ \\

  \paragraph{Array:}
  \[\type = \typearray{\type_a}{\idx}\]
  $\EraseT{\idxsubst{\typearray{\type_a}{\idx}}}$ \\
  $= \EraseT{\typearray{\idxsubst{\type_a}}{\idxsubst{\idx}}}$ \\
  $= \idxsubst{\idx}$ \\
  $= \idxsubst{\EraseT{\typearray{\type_a}{\idx}}}$ \\

\end{sproof}

\begin{lemma}%[Index-into-term substitution commutes with erasure]
  \label{IESubstErase}
  $\EraseTrm{\idxsubst{\term}} = \idxsubst{\EraseTrm{\term}}$
\end{lemma}
\begin{sproof}[This is straightforward induction on $\term$.]
  We use induction on $\term$.
  We elide the cases where $\var$ does not appear free in $\term$.

  \paragraph{Term abstraction:}
  \[\term = \lam{\sequence{\notevar{\var_i}{\type}}}{\expr}\]
  $\EraseTrm{\idxsubst{\lam{\sequence{\notevar{\var_i}{\type}}}{\expr}}}$ \\
  $= \EraseA{\idxsubst{\lam{\sequence{\notevar{\var_i}{\type}}}{\expr}}}$ \\
  $= \EraseA{\lam{\sequence{\notevar{\var_i}{\type}}}{\idxsubst{\expr}}}$ \\
  $= \lam{\sequence{\var_i}}{\EraseE{\idxsubst{\expr}}}$ \\
  $= \lam{\sequence{\var_i}}{\idxsubst{\EraseE{\expr}}}$ \\
  $= \idxsubst{\lam{\sequence{\var_i}}{\EraseE{\expr}}}$ \\
  $= \idxsubst{\EraseA{\lam{\sequence{\var_i}}{\expr}}}$ \\
  $= \idxsubst{\EraseTrm{\lam{\sequence{\var_i}}{\expr}}}$ \\

  \paragraph{Type abstraction:}
  \[\term = \tlam{\sequence{\notevar{\var_i}{\kind}}}{\val}\]
  Note: $\var$ is an index variable, while the $\sequence{\var_i}$ are type variables, so they do not shadow.
  $\EraseTrm{\idxsubst{\tlam{\sequence{\notevar{\var_i}{\kind}}}{\val}}}$ \\
  $= \EraseA{\idxsubst{\tlam{\sequence{\notevar{\var_i}{\kind}}}{\val}}}$ \\
  $= \EraseA{\tlam{\sequence{\notevar{\var_i}{\kind}}}{\idxsubst{\val}}}$ \\
  $= \ilam{\sequence{\var_i}}{\EraseE{\idxsubst{\val}}}$ \\
  $= \ilam{\sequence{\var_i}}{\idxsubst{\EraseE{\val}}}$ \\
  $= \idxsubst{\ilam{\sequence{\var_i}}{\EraseE{\val}}}$ \\
  $= \idxsubst{\EraseA{\tlam{\sequence{\notevar{\var_i}{\kind}}}{\val}}}$ \\
  $= \idxsubst{\EraseTrm{\tlam{\sequence{\notevar{\var_i}{\kind}}}{\val}}}$ \\

  \paragraph{Index abstraction:}
  \[\term = \ilam{\sequence{\notevar{\var_i}{\sort}}}{\val}\text{, where } \var \not\in \sequence{\var_i}\]
  $\EraseTrm{\idxsubst{\ilam{\sequence{\notevar{\var_i}{\sort}}}{\val}}}$ \\
  $= \EraseA{\idxsubst{\ilam{\sequence{\notevar{\var_i}{\sort}}}{\val}}}$ \\
  $= \EraseA{\ilam{\sequence{\notevar{\var_i}{\sort}}}{\idxsubst{\val}}}$ \\
  $= \ilam{\sequence{\var_i}}{\EraseE{\idxsubst{\val}}}$ \\
  $= \ilam{\sequence{\var_i}}{\idxsubst{\EraseE{\val}}}$ \\
  $= \idxsubst{\ilam{\sequence{\var_i}}{\EraseE{\val}}}$ \\
  $= \idxsubst{\EraseA{\ilam{\sequence{\notevar{\var_i}{\sort}}}{\val}}}$ \\
  $= \idxsubst{\EraseTrm{\ilam{\sequence{\notevar{\var_i}{\sort}}}{\val}}}$ \\

  \paragraph{Box:}
  \[\term = \dsum{\sequence{\idx}}{\expr_s}{\type}\]
  $\EraseTrm{\idxsubst{\dsum{\sequence{\idx}}{\expr_s}{\type}}}$ \\
  $= \EraseA{\idxsubst{\dsum{\sequence{\idx}}{\expr_s}{\type}}}$ \\
  $= \EraseA{\dsum{\sequence{\idxsubst{\idx}}}{\idxsubst{\expr_s}}{\idxsubst{\type}}}$ \\
  $= \dsum{\sequence{\idxsubst{\idx}}}{\EraseE{\idxsubst{\expr_s}}}{}$ \\
  $= \dsum{\sequence{\idxsubst{\idx}}}{\idxsubst{\EraseE{\expr_s}}}{}$ \\
  $= \idxsubst{\dsum{\sequence{\idx}}{\EraseE{\expr_s}}{}}$ \\
  $= \idxsubst{\EraseA{\dsum{\sequence{\idx}}{\expr_s}{\type}}}$ \\
  $= \idxsubst{\EraseTrm{\dsum{\sequence{\idx}}{\expr_s}{\type}}}$ \\

  \paragraph{Array:}
  \[\term = \arrlit{\sequence{\atom}}{\sequence{\nat}}\]
  $\EraseTrm{\idxsubst{\arrlit{\sequence{\atom}}{\sequence{\nat}}}}$ \\
  $= \EraseE{\idxsubst{\arrlit{\sequence{\atom}}{\sequence{\nat}}}}$ \\
  $= \EraseE{\arrlit{\sequence{\idxsubst\atom}}{\sequence{\nat}}}$ \\
  $= \arrlit{\sequence{\EraseA{\idxsubst{\atom}}}}{\sequence{\nat}}$ \\
  $= \arrlit{\sequence{\idxsubst{\EraseA{\atom}}}}{\sequence{\nat}}$ \\
  $= \idxsubst{\arrlit{\sequence{\EraseA{\atom}}}{\sequence{\nat}}}$ \\
  $= \idxsubst{\EraseE{\arrlit{\sequence{\atom}}{\sequence{\nat}}}}$ \\
  $= \idxsubst{\EraseTrm{\arrlit{\sequence{\atom}}{\sequence{\nat}}}}$ \\

  \paragraph{Frame:}
  \[\term = \notedfrm[\type_r]{\sequence{\expr_c}}{\sequence{\nat}}\]
  $\EraseTrm{\idxsubst{\notedfrm[\type_r]{\sequence{\expr_c}}{\sequence{\nat}}}}$ \\
  $= \EraseE{\idxsubst{\notedfrm[\type_r]{\sequence{\expr_c}}{\sequence{\nat}}}}$ \\
  $= \EraseE{\notedfrm[\idxsubst{\type_r}]{\sequence{\idxsubst{\expr_c}}}{\sequence{\nat}}}$ \\
  $= \frm{\sequence{\EraseE{\idxsubst{\expr_c}}}}{\EraseT{\idxsubst{\type_r}}}$ \\
  $= \frm{\sequence{\idxsubst{\EraseE{\expr_c}}}}{\EraseT{\idxsubst{\type_r}}}$%
  , by the induction hypothesis \\
  $= \frm{\sequence{\idxsubst{\EraseE{\expr_c}}}}{\idxsubst{\EraseT{\type_r}}}$%
  , by Lemma \ref{ITSubstErase} \\
  $= \idxsubst{\frm{\sequence{\EraseE{\expr_c}}}{\EraseT{\type_r}}}$ \\
  $= \idxsubst{\EraseE{\notedfrm[\type_r]{\sequence{\expr_c}}{\sequence{\nat}}}}$ \\
  $= \idxsubst{\EraseTrm{\notedfrm[\type_r]{\sequence{\expr_c}}{\sequence{\nat}}}}$ \\

  \paragraph{Application:}
  \[\term =
  \notedapp[\type_r]{\annotate
    [\typearray{\typefun{\sequence{\type_i}}{\type_o}}{\idx_f}]
    {\expr_f}}{\sequence{\expr_a}}\]
  $\EraseTrm{\idxsubst{\notedapp[\type_r]{\annotate
        [\typearray{\typefun{\sequence{\type_i}}{\type_o}}{\idx_f}]
        {\expr_f}}{\sequence{\expr_a}}}}$ \\
  $= \EraseE{\idxsubst{\notedapp[\type_r]{\annotate
        [\typearray{\typefun{\sequence{\type_i}}{\type_o}}{\idx_f}]
        {\expr_f}}{\sequence{\expr_a}}}}$ \\
  $= \EraseE{\notedapp[\idxsubst{\type_r}]{\annotate
      [\typearray{\typefun{\sequence{\idxsubst{\type_i}}}{\idxsubst{\type_o}}}{\idxsubst{\idx_f}}]
      {\idxsubst{\expr_f}}}{\sequence{\idxsubst{\expr_a}}}}$ \\
  $= \app{\EraseE{\idxsubst{\expr_f}}}{\sequence
    {\ttsexp{\EraseE{\idxsubst{\expr_a}}}{\EraseT{\idxsubst{\type_i}}}}\;\EraseT{\idxsubst{\type_r}}}$ \\
  $= \app{\idxsubst{\EraseE{\expr_f}}}{\sequence
    {\ttsexp{\idxsubst{\EraseE{\expr_a}}}{\EraseT{\idxsubst{\type_i}}}}\;\EraseT{\idxsubst{\type_r}}}$%
  , by the induction hypothesis \\
  $= \app{\idxsubst{\EraseE{\expr_f}}}{\sequence
    {\ttsexp{\idxsubst{\EraseE{\expr_a}}}{\idxsubst{\EraseT{\type_i}}}}\;\idxsubst{\EraseT{\type_r}}}$%
  , by Lemma \ref{TTSubstErase} \\
  $= \idxsubst{\app{\EraseE{\expr_f}}{\sequence
      {\ttsexp{\EraseE{\expr_a}}{\EraseT{\type_i}}}\;\EraseT{\type_r}}}$ \\
  $= \idxsubst{\EraseE{\notedapp[\type_r]{\annotate
        [\typearray{\typefun{\sequence{\type_i}}{\type_o}}{\idx_f}]
        {\expr_f}}{\sequence{\expr_a}}}}$ \\
  $= \idxsubst{\EraseTrm{\notedapp[\type_r]{\annotate
        [\typearray{\typefun{\sequence{\type_i}}{\type_o}}{\idx_f}]
        {\expr_f}}{\sequence{\expr_a}}}}$ \\

  \paragraph{Type application:}
  \[\term = \notedtapp[\type_r]{\expr_f}{\sequence{\type_a}}\]
  $\EraseTrm{\idxsubst{\notedtapp[\type_r]{\expr_f}{\sequence{\type_a}}}}$ \\
  $= \EraseE{\idxsubst{\notedtapp[\type_r]{\expr_f}{\sequence{\type_a}}}}$ \\
  $= \EraseE{\notedtapp[\idxsubst{\type_r}]{\idxsubst{\expr_f}}{\sequence{\idxsubst{\type_a}}}}$ \\
  $= \iapp{\EraseE{\idxsubst{\expr_f}}}{\sequence{\EraseT{\idxsubst{\type_a}}}\;\EraseT{\idxsubst{\type_r}}}$ \\
  $= \iapp{\idxsubst{\EraseE{\expr_f}}}{\sequence{\EraseT{\idxsubst{\type_a}}}\;\EraseT{\idxsubst{\type_r}}}$%
  , by the induction hypothesis \\
  $= \iapp{\idxsubst{\EraseE{\expr_f}}}{\sequence{\idxsubst{\EraseT{\type_a}}}\;\idxsubst{\EraseT{\type_r}}}$%
  , by Lemma \ref{TTSubstErase} \\
  $= \idxsubst{\iapp{\EraseE{\expr_f}}{\sequence{\EraseT{\type_a}}\;\EraseT{\type_r}}}$ \\
  $= \idxsubst{\EraseE{\notedtapp[\type_r]{\expr_f}{\sequence{\type_a}}}}$ \\
  $= \idxsubst{\EraseTrm{\notedtapp[\type_r]{\expr_f}{\sequence{\type_a}}}}$ \\

  \paragraph{Index application:}
  \[\term = \notediapp[\type_r]{\expr_f}{\sequence{\idx_a}}\]
  $\EraseTrm{\idxsubst{\notediapp[\type_r]{\expr_f}{\sequence{\idx_a}}}}$ \\
  $= \EraseE{\idxsubst{\notediapp[\type_r]{\expr_f}{\sequence{\idx_a}}}}$ \\
  $= \EraseE{\notediapp[\idxsubst{\type_r}]{\idxsubst{\expr_f}}{\sequence{\idxsubst{\idx_a}}}}$ \\
  $= \iapp{\EraseE{\idxsubst{\expr_f}}}{\sequence{\idxsubst{\idx_a}}\;\EraseT{\idxsubst{\type_r}}}$ \\
  $= \iapp{\idxsubst{\EraseE{\expr_f}}}{\sequence{\idxsubst{\idx_a}}\;\EraseT{\idxsubst{\type_r}}}$%
  , by the induction hypothesis \\
  $= \iapp{\idxsubst{\EraseE{\expr_f}}}{\sequence{\idxsubst{\idx_a}}\;\idxsubst{\EraseT{\type_r}}}$%
  , by Lemma \ref{TTSubstErase} \\
  $= \idxsubst{\iapp{\EraseE{\expr_f}}{\sequence{\idx_a}\;\EraseT{\type_r}}}$ \\
  $= \idxsubst{\EraseE{\notediapp[\type_r]{\expr_f}{\sequence{\idx_a}}}}$ \\
  $= \idxsubst{\EraseTrm{\notediapp[\type_r]{\expr_f}{\sequence{\idx_a}}}}$ \\

  \paragraph{Unboxing:}
  \[\term = \dproj{\sequence{\var_i}}{\var_e}{\expr_s}{\expr_b^{\type_b}}\text{, where }\var \not\in \sequence{\var_i}\]
  $\EraseTrm{\idxsubst{\dproj{\sequence{\var_i}}{\var_e}{\expr_s}{\expr_b^{\type_b}}}}$ \\
  $= \EraseE{\idxsubst{\dproj{\sequence{\var_i}}{\var_e}{\expr_s}{\expr_b^{\type_b}}}}$ \\
  $= \EraseE{\dproj{\sequence{\var_i}}{\var_e}{\idxsubst{\expr_s}}{\idxsubst{\expr_b}^{\idxsubst{\type_b}}}}$ \\
  $= \dproj{\sequence{\var_i}}{\var_e}{\EraseE{\idxsubst{\expr_s}}}
  {\EraseE{\idxsubst{\expr_b}} \; \EraseT{\idxsubst{\type_b}}}$ \\
  $= \dproj{\sequence{\var_i}}{\var_e}{\EraseE{\idxsubst{\expr_s}}}
  {\EraseE{\idxsubst{\expr_b}} \; \idxsubst{\EraseT{\type_b}}}$ \\
  $= \dproj{\sequence{\var_i}}{\var_e}{\idxsubst{\EraseE{\expr_s}}}
  {\idxsubst{\EraseE{\expr_b}} \; \idxsubst{\EraseT{\type_b}}}$ \\
  $= \idxsubst{\dproj{\sequence{\var_i}}{\var_e}{\EraseE{\expr_s}}{\EraseE{\expr_b} \; \EraseT{\type_b}}}$ \\
  $= \idxsubst{\EraseE{\dproj{\sequence{\var_i}}{\var_e}{\expr_s}{\expr_b^{\type_b}}}}$ \\
  $= \idxsubst{\EraseTrm{\dproj{\sequence{\var_i}}{\var_e}{\expr_s}{\expr_b^{\type_b}}}}$ \\

  \paragraph{Unboxing, with shadowed variable:}
  \[\term = \dproj{\sequence{\var_i}}{\var_e}{\expr_s}{\expr_b^{\type_b}}\text{, where }\var \in \sequence{\var_i}\]
  $\EraseTrm{\idxsubst{\dproj{\sequence{\var_i}}{\var_e}{\expr_s}{\expr_b^{\type_b}}}}$ \\
  $= \EraseE{\idxsubst{\dproj{\sequence{\var_i}}{\var_e}{\expr_s}{\expr_b^{\type_b}}}}$ \\
  $= \EraseE{\dproj{\sequence{\var_i}}{\var_e}{\idxsubst{\expr_s}}{\expr_b^{\type_b}}}$ \\
  $= \dproj{\sequence{\var_i}}{\var_e}{\EraseE{\idxsubst{\expr_s}}}{\EraseE{\expr_b} \; \EraseT{\type_b}}$ \\
  $= \dproj{\sequence{\var_i}}{\var_e}{\idxsubst{\EraseE{\expr_s}}}{\EraseE{\expr_b} \; \EraseT{\type_b}}$ \\
  $= \idxsubst{\dproj{\sequence{\var_i}}{\var_e}{\EraseE{\expr_s}}{\EraseE{\expr_b} \; \EraseT{\type_b}}}$ \\
  $= \idxsubst{\EraseE{\dproj{\sequence{\var_i}}{\var_e}{\expr_s}{\expr_b^{\type_b}}}}$ \\
  $= \idxsubst{\EraseTrm{\dproj{\sequence{\var_i}}{\var_e}{\expr_s}{\expr_b^{\type_b}}}}$ \\

\end{sproof}

\begin{lemma}[Values erase to values]
  \label{ValEraseToVal}
  For any well-typed term $\term$,
  \begin{itemize}
  \item If $\term$ has the form $\atval$, then $\EraseTrm{\term}$ has the form $\eatvalue$
  \item If $\term$ has the form $\val$, then $\EraseTrm{\term}$ has the form $\evalue$
  \end{itemize}
  % $\EraseE{\val}$ has the form $\evalue$.
\end{lemma}
\begin{sproof}[We use induction on $\term$.
  The only nontrivial cases are {\tt box} and {\tt array} forms,
  which may be values or may contain incomplete computation.
  The contents of a {\tt box} value
  must itself be a value,
  which the induction hypothesis implies will erase to a value.
  Similarly, an {\tt array} value contains only atomic values,
  which erase to atomic values.]
  We use induction on $\term$.
  The result is trivial for all atom cases except for boxes,
  so we have only two cases left to consider.

  \paragraph{Box:}
  \[\term = \dsum{\sequence{\idx}}{\val}{\type}\]
  According to the grammar of Remora,
  $\val$ must have the form
  $\arrlit{\sequence{\atval}}{\sequence{\nat}}$ \\
  Then $\EraseTrm{\term}$
  $= \EraseA{\dsum{\sequence{\idx}}{\val}{\type}}$
  $= \dsum{\sequence{\idx}}{\EraseE{\val}}{}$.
  By the induction hypothesis,
  $\EraseE{\val}$ has the form $\evalue$,
  so $\EraseTrm{\term}$
  has the form $\dsum{\sequence{\idx}}{\evalue}{}$,
  which is a valid $\eatvalue$ form.

  \paragraph{Array literal:}
  \[\term = \arrlit{\sequence{\atval}}{\sequence{\nat}}\]
  The induction hypothesis implies that
  for each $\atval_i \in \sequence{\atval}$,
  $\EraseA{\atval_i}$ produces an erased atomic value $\eatvalue_i$.
  Therefore the erased term $\EraseTrm{\term}$
  $=$
  $\EraseE{\arrlit{\sequence{\atval}}{\sequence{\nat}}}$
  $=$
  {\tt (array ($\nat$ $\sequence{}$) $\EraseA{\atval}$ $\sequence{}$)}
  must have the form
  $\arrlit{\sequence{\eatvalue}}{\sequence{\nat}}$.
\end{sproof}

\begin{lemma}[Lockstep]
  \label{Lockstep}
  For any well-typed $\expr$,
  one of the following holds:
  \begin{itemize}
  \item{$\expr$ has the form $\val$, and $\EraseE{\expr}$ has the form $\evalue$}
  \item{$\expr \mapsto \expr'$, and $\EraseE{\expr} \mapsto \EraseE{\expr'}$}
  \item{$\expr \not \mapsto$, and $\EraseE{\expr} \not \mapsto$}
  \end{itemize}
\end{lemma}
\begin{sproof}[We prove this by induction on $\expr$.
  
  We rely on Lemma \ref{ErasureInContext} (erasure in context)
  when $\expr$ is a redex $\expr_r$ within an evaluation context $\ctxt$ other than $\hole$.
  If $\expr_r \not\mapsto$ because we have a mis-applied primitive operator,
  then the same is true for $\EraseE{\expr_r}$,
  so $\EraseE{\expr}$ is also an evaluation context around a mis-applied primitive operator.
  Otherwise, $\expr_r \mapsto \expr_r'$,
  and the induction hypothesis implies that $\EraseE{\expr_r} \mapsto \EraseE{\expr_r'}$.
  So $\EraseE{\expr} \mapsto \EraseE{\expr'}$,
  the erased context filled with $\expr_r'$.
  
  Values are handled by Lemma \ref{ValEraseToVal}.
  For the remaining cases---redexes---%
  straightforward symbol pushing shows that
  erased Remora's reduction rules follow those of Remora.]
  We prove this by induction on $\expr$.
  Note that if $\expr$ is not itself a redex or a value form,
  then the progress lemma implies that it must be
  an evaluation context filled with a redex.
  \paragraph{Value:}
  This case is exactly the expression case of Lemma \ref{ValEraseToVal}.

  \paragraph{Redex within nontrivial evaluation context:}
  \[\expr = \fillc{\ctxt}{\expr_r}\text{, where }\expr_r \mapsto \expr_r'\]
  Then $\expr \mapsto \expr' = \fillc{\ctxt}{\expr_r'}$.
  By Lemma \ref{ErasureInContext} (erasure in context),
  $\EraseE{\fillc{\ctxt}{\expr_r}} = \fillc{\EraseC{\ctxt}}{\EraseE{\expr_r}}$.
  The induction hypothesis implies that $\EraseE{\expr_r} \mapsto \EraseE{\expr_r'}$,
  so the full erased expression $\fillc{\EraseC{\ctxt}}{\EraseE{\expr_r}}$
  $ \mapsto \fillc{\EraseC{\ctxt}}{\EraseE{\expr_r'}}$.
  Erasure in context gives us
  $\fillc{\EraseC{\ctxt}}{\EraseE{\expr_r'}} = \EraseE{\fillc{\ctxt}{\expr_r'}}$.
  Therefore $\EraseE{\fillc{\ctxt}{\expr_r}} \mapsto \EraseE{\fillc{\ctxt}{\expr_r'}}$.

  \paragraph{Lift redex:}
  % Pulling these out as macro defns because linebreaking for code readability
  % is not really an option inside an alltt block
  \newcommand{\liftfnannot}
  {\typearray
    {\typefun
      {\sequence{\typearray{\type_i}{\idxshape{\sequence{\nat_i}}}}}
      {\typearray{\type_o}{\idx_o}}}
    {\idxshape{\sequence{\nat_f}}}}
  \newcommand{\liftfn}
  {\notedarrlit[\liftfnannot]{\sequence{\atval_f}}{\sequence{\nat_f}}}
  \newcommand{\liftargannot}
  {\typearray
    {\type_i}
    {\idxshape{\sequence{\nat_a}\;\sequence{\nat_i}}}}
  \newcommand{\liftarg}
  {\notedarrlit[\liftargannot]{\sequence{\atval_a}}{\sequence{\nat_a}\,\sequence{\nat_i}}}
  \newcommand{\liftresultannot}
  {\typearray
    {\type_o}
    {\idxshape{\sequence{\nat_p}\;\sequence{\nat_o}}}}
  \newcommand{\liftedfnannot}
  {\typearray
    {\typefun
      {\sequence{\typearray{\type_i}{\idxshape{\sequence{\nat_i}}}}}
      {\typearray{\type_o}{\idx_o}}}
    {\idxshape{\sequence{\nat_p}}}}
  \newcommand{\liftedfncells}
  {\mathit{Concat}\llb
    \mathit{Rep}_{\nat_{\mathit{fe}}}\llb
    \mathit{Split}_1\llb
    \sequence{\atval_f}\rrb\rrb\rrb}
  \newcommand{\liftedfn}
  {\notedarrlit[\liftedfnannot]
    {\liftedfncells}
    {\sequence{\nat_p}}}
  \newcommand{\liftedargannot}
  {\typearray
    {\type_i}
    {\idxshape{\sequence{\nat_p}\;\sequence{\nat_i}}}}
  \newcommand{\liftedarg}
  {\notedarrlit[\liftedargannot]
    {\mathit{Concat}\llb
      \mathit{Rep}_{\nat_{\mathit{ae}}}\llb
      \mathit{Split}_{\nat_{\mathit{ac}}}\llb
      \sequence{\atval_a}\rrb\rrb\rrb}
    {\sequence{\nat_p}\,\sequence{\nat_i}}}
  \newcommand{\erasedliftfn}
  {\arrlit{\sequence{\EraseA{\atval_f}}}{\sequence{\nat_f}}}
  \newcommand{\erasedliftarg}
  {\arrlit{\sequence{\EraseA{\atval_a}}}{\sequence{\nat_a}\,\sequence{\nat_i}}}
  \newcommand{\erasedliftcell}{\idxshape{\sequence{\nat_i}}}
  \newcommand{\erasedliftresult}{\idxshape{\sequence{\nat_p}\;\sequence{\nat_o}}}
  % \newcommand{\erasedliftresult}{\EraseT{\liftresultannot}}
  \newcommand{\erasedliftedfn}
  {\arrlit
    {\mathit{Concat}\llb
      \mathit{Rep}_{\nat_{\mathit{fe}}}\llb
      \mathit{Split}_1\llb
      \sequence{\EraseA{\atval_f}}\rrb\rrb\rrb}
    {\sequence{\nat_p}}}
  \newcommand{\erasedliftedarg}
  {\arrlit
    {\mathit{Concat}\llb
      \mathit{Rep}_{\nat_{\mathit{ae}}}\llb
      \mathit{Split}_{\nat_{\mathit{ac}}}\llb
      \sequence{\EraseA{\atval_a}}\rrb\rrb\rrb}
    {\sequence{\nat_p}\,\sequence{\nat_i}}}
\begin{alltt}
\(\expr = \)(\(\liftfn\)\\
\phantom{\(\expr = \) }\(\liftarg\)\\
\phantom{\(\expr = \) }\(\cdots\))\(\sp{\liftresultannot}\)\\
\(\mapsto\sb{\mathit{lift}}\)\\
\(\expr' = \)((array ()\\
\phantom{\(\expr' = \) \ }\(\liftedfncells\))\(\sp{\liftedfnannot}\)\\
\phantom{\(\expr' = \) }\(\liftedarg\)\\
\phantom{\(\expr' = \) }\(\cdots\))\(\sp{\liftresultannot}\)
\end{alltt}
Then
\begin{alltt}
\(\EraseE{\expr} = \)(\(\erasedliftfn\)\\
\phantom{\(\EraseE{\expr} = \) }(\(\erasedliftarg\) \(\erasedliftcell\))\\
\phantom{\(\EraseE{\expr} = \) }\(\cdots\) \(\erasedliftresult\))
\end{alltt}
This is a $\mathit{lift}$ redex in Erased Remora, and it steps to
\begin{alltt}
\(\EraseE{\expr'} = \)(\(\erasedliftedfn\)\\
\phantom{\(\EraseE{\expr'} = \) }(\(\erasedliftedarg\) \(\erasedliftcell\))\\
\phantom{\(\EraseE{\expr'} = \) }\(\cdots\) \(\erasedliftresult\))
\end{alltt}
That is, $\EraseE{\expr} \mapsto_{\mathit{lift}} \EraseE{\expr'}$.

  \paragraph{Map redex:}
  \newcommand{\mapfn}{\liftfn}
  \newcommand{\mapargannot}
  {\typearray
    {\type_i}
    {\idxshape{\sequence{\nat_f}\;\sequence{\nat_i}}}}
  \newcommand{\maparg}
  {\notedarrlit[\mapargannot]{\sequence{\atval_a}}{\sequence{\nat_f}\,\sequence{\nat_i}}}
  \newcommand{\mapresultannot}
  {\typearray
    {\type_o}
    {\idxshape{\sequence{\nat_f}\;\sequence{\nat_o}}}}
  \newcommand{\mapfncell}
  {\notedarrlit[{\typearray
      {\typefun
        {\sequence{\typearray{\type_i}{\idxshape{\sequence{\nat_i}}}}}
        {\typearray{\type_o}{\idx_o}}}
      {\idxshape{}}}]
    {\atval_f}{}}
  \newcommand{\mapargcell}
  {\notedarrlit[{\typearray{\type_i}{\idxshape{\sequence{\nat_i}}}}]
    {\sequence{\atval_c}}{}}
  \newcommand{\mappedargannot}
  {\typearray
    {\type_i}
    {\idxshape{\sequence{\nat_i}}}}
  \newcommand{\resultcellannot}
  {\typearray
    {\type_o}
    {\idxshape{\sequence{\nat_o}}}}
  \newcommand{\erasedmapfn}{\erasedliftfn}
  \newcommand{\erasedmaparg}
  {\arrlit{\sequence{\EraseA{\atval_a}}}{\sequence{\nat_f}\,\sequence{\nat_i}}}
  \newcommand{\erasedmapcell}{\idxshape{\sequence{\nat_i}}}
  \newcommand{\erasedmapresult}{\erasedliftresult}
  \newcommand{\erasedmapfncell}{\arrlit{\EraseA{\atval_f}}{}}
  \newcommand{\erasedmapargcell}{\arrlit{\sequence{\EraseA{\atval_c}}}{}}
\begin{alltt}
\(\expr = \)(\(\mapfn\)\\
\phantom{\(\expr = \) }\(\maparg\)\\
\phantom{\(\expr = \) }\(\cdots\))\(\sp{\mapresultannot}\)\\
\(\mapsto\sb{\mathit{map}}\)\\
\(\expr' = \)(frame (\(\sequence{\nat\sb{f}}\))\\
\phantom{\(\expr' = \) }(\(\mapfncell\)\\
\phantom{\(\expr' = \) \ }\(\mapargcell\)\\
\phantom{\(\expr' = \) \ }\(\cdots\))\(\sp{\resultcellannot}\)\\
\phantom{\(\expr' = \) }\(\cdots\))\(\sp{\mapresultannot}\)
\end{alltt}
The argument cells' atoms
$\parens{\sequence{\parens{\sequence{\atval_c}}}}$
are given by
$\mathit{Transpose} \llb \sequence{\mathit{Split}_{\nat_c} \llb \sequence{\atval_a} \rrb} \rrb$,
where
each $\nat_c$ is computed as the product of
the corresponding position's expected argument dimensions $\sequence{\nat_i}$,
as in Figure \ref{fig:DynamicSemantics}.
We then consider the erased form of $\expr$:
\begin{alltt}
\(\EraseE{\expr} = \)(\(\erasedmapfn\)
         (\(\erasedmaparg\) \(\erasedmapcell\))
         \(\cdots\) \(\erasedliftresult\))
\end{alltt}
This is a $\mathit{map}$ redex in Erased Remora.
The argument atoms
$\parens{\sequence{\parens{\sequence{\EraseA{\atval_a}}}}}$
are split up into cells in the same way, with
$\parens{\sequence{\parens{\sequence{\eatvalue_c}}}}
= \parens{\sequence{\parens{\sequence{\EraseA{\atval_c}}}}}$
as the result of
$\mathit{Transpose} \llb \sequence{\mathit{Split}_{\nat_c} \llb \sequence{\EraseA{\atval_a}} \rrb} \rrb$.
So $\EraseE{\expr}$ steps to
\begin{samepage}
\begin{alltt}
\(\EraseE{\expr'} = \)(frame (\(\sequence{\nat\sb{f}}\))\\
\phantom{\(\EraseE{\expr'} = \) }(\(\erasedmapfncell\)\\
\phantom{\(\EraseE{\expr'} = \) \ }(\(\erasedmapargcell\) \(\erasedmapcell\))\\
\phantom{\(\EraseE{\expr'} = \) \ }\(\cdots\))\\
\phantom{\(\EraseE{\expr'} = \) }\(\cdots\))
\end{alltt}
\end{samepage}

  \paragraph{Beta redex:}
  Note that $\type_I$ here must be the array type
  $\typearray{\type_i}{\idxshape{\sequence{\nat_a}}}$.
  \newcommand{\betafnannot}
  {\typearray
    {\typefun
      {\sequence{\type_I}}
      {\type_O}}
    {\idxshape{}}}
  \newcommand{\betafn}
  {\notedarrlit[\betafnannot]
    {\lam{\sequence{\notevar{\var}{\type_i}}}{\expr_b}}
    {}}
  \newcommand{\erasedbetafn}
  {\arrlit
    {\lam{\sequence{\var}}{\EraseE{\expr_b}}}
    {}}
  \newcommand{\betaarg}{\notedarrlit[\type_I]{\sequence{\atval_a}}{\sequence{\nat_a}}}
  \newcommand{\erasedbetaarg}{\arrlit{\sequence{\EraseA{\atval_a}}}{\sequence{\nat_a}}}
  \newcommand{\erasedbetacell}{\idxshape{\sequence{\nat_a}}}
\begin{alltt}
\(\expr = \)(\(\betafn\)\\
\phantom{\(\expr = \) }\(\betaarg\)\\
\phantom{\(\expr = \) }\(\sequence{}\))\(\sp{\type\sb{O}}\)\\
\(\mapsto\sb{\beta}\)\\
\(\expr' = \seqsubst{\expr\sb{b}}{\var}{\betaarg}\)
\end{alltt}
Then
\begin{alltt}
\(\EraseE{\expr} = \)(\(\erasedbetafn\)\\
\phantom{\(\EraseE{\expr} = \) }(\(\erasedbetaarg\) \(\erasedbetacell\))\\
\phantom{\(\EraseE{\expr} = \) }\(\sequence{}\))\(\sp{\type\sb{O}}\)\\
\(\mapsto\sb{\beta}\) \(\seqsubst{\EraseE{\expr\sb{b}}}{\var}{\erasedbetaarg}\)
\end{alltt}
By Lemma \ref{EESubstErase} (substitution commutes with erasure),
this result term is equal to $\EraseE{\expr'}$.

%  \paragraph{Delta redex:} \pfstub

  \paragraph{I-Beta redex:}
  \newcommand{\ibetafnannot}
  {\typearray
    {\typedprod
      {\sequence{\notevar{\var}{\sort}}}
      {\type_b}}
    {\idxshape{\sequence{\nat_f}}}}
  \newcommand{\ibetafn}
  {\notedarrlit[\ibetafnannot]
    {\sequence{\ilam{\notevar{\var}{\sort}}{\expr_b}}}
    {\sequence{\nat_f}}}
  \newcommand{\ibetaresult}
  {\notedfrm[\type_R]
    {\sequence{\seqsubst{\expr_b}{\var}{\idx_a}}}
    {\sequence{\nat_f}}}
  \newcommand{\erasedibetafn}
  {\arrlit
    {\sequence{\ilam{\var}{\EraseE{\expr_b}}}}
    {\sequence{\nat_f}}}
  \newcommand{\erasedibetaresult}
  {\notedfrm[\type_R]
    {\sequence{\seqsubst{\EraseE{\expr_b}}{\var}{\idx_a}}}
    {\sequence{\nat_f}}}
\begin{alltt}
\(\expr = \)(i-app \(\ibetafn\)\\
\phantom{\(\expr = \) }\(\idx\sb{a}\) ...)\(\sp{\type\sb{R}}\)\\
\(\mapsto\sb{i\beta}\)\\
\(\expr' = \ibetaresult\)
\end{alltt}
Then
\begin{alltt}
\(\EraseE{\expr} = \)(i-app \(\erasedibetafn\) \(\idx\sb{a}\) ... \(\EraseT{\type\sb{R}}\))\\
\(\mapsto\sb{i\beta} \erasedibetaresult = \EraseE{\expr'}\)
\end{alltt}

  \paragraph{T-Beta redex:}
  \newcommand{\tbetafnannot}
  {\typearray
    {\typeuniv
      {\sequence{\notevar{\var}{\kind}}}
      {\type_b}}
    {\idxshape{\sequence{\nat_f}}}}
  \newcommand{\tbetafn}
  {\notedarrlit[\tbetafnannot]
    {\sequence{\tlam{\notevar{\var}{\kind}}{\expr_b}}}
    {\sequence{\nat_f}}}
  \newcommand{\tbetaresult}
  {\notedfrm[\type_R]
    {\sequence{\seqsubst{\expr_b}{\var}{\type_a}}}
    {\sequence{\nat_f}}}
  \newcommand{\erasedtbetafn}
  {\arrlit
    {\sequence{\ilam{\var}{\EraseE{\expr_b}}}}
    {\sequence{\nat_f}}}
  \newcommand{\erasedtbetaresult}
  {\notedfrm[\type_R]
    {\sequence{\seqsubst{\EraseE{\expr_b}}{\var}{\type_a}}}
    {\sequence{\nat_f}}}
\begin{alltt}
\(\expr = \)(t-app (array (\(\sequence{\nat\sb{f}}\))\\
\phantom{\(\expr = \)(t-app (array }\(\sequence{\tlam{\notevar{\var}{\kind}}{\expr\sb{b}}}\))\(\sp{\tbetafnannot}\)\\
\phantom{\(\expr = \) }\(\type\sb{a}\) ...)\(\sp{\type\sb{R}}\)\\
\(\mapsto\sb{t\beta}\)\\
\(\expr' = \tbetaresult\)
\end{alltt}
Recall that type abstraction and application erase to index abstraction and application.
So in the erased language, we have:
\begin{alltt}
\(\EraseE{\expr} = \)(i-app \(\erasedtbetafn\) \(\type\sb{a}\) ... \(\EraseT{\type\sb{R}}\))\\
\(\mapsto\sb{i\beta} \erasedtbetaresult = \EraseE{\expr'}\)
\end{alltt}

  \paragraph{Unbox redex:}
  \newcommand{\unboxvars}{\sequence{\var_i}\text{ }\var_e}
  \newcommand{\unboxval}{\annotate%[\typearray{\type_s}{\idxshape{\sequence{\nat_s}}}]
    {\arrlit{\sequence{\dsum{\sequence{\idx_s}}{\val_s}{\type_s}}}{\sequence{\nat_s}}}}
  \newcommand{\unboxbody}{\annotate[\type_B]{\expr_b}}
  \newcommand{\unboxresult}
  {\notedfrm[\type_R]
    {\sequence{\multisubst{\expr_b}{\sequence{\singlesubst{\var_i}{\idx_s},},
          \singlesubst{\var_e}{\val_s}}}}
    {\sequence{\nat_s}}}
  \newcommand{\erasedunboxval}
  {\arrlit{\sequence{\dsum{\sequence{\idx_s}}{\EraseE{\val_s}}{}}}{\sequence{\nat_s}}}
  \newcommand{\erasedunboxresult}
  {\frm
    {\sequence{\multisubst{\EraseE{\expr_b}}{\sequence{\singlesubst{\var_i}{\idx_s},},
          \singlesubst{\var_e}{\EraseE{\val_s}}}}}
    {\idxappend{\idxshape{\sequence{\nat_s}} \; \EraseT{\type_B}}}}
  \newcommand{\rewrittenunboxresult}
  {\frm
    {\sequence{\EraseE{\multisubst{\expr_b}{\sequence{\singlesubst{\var_i}{\idx_s},},
            \singlesubst{\var_e}{\val_s}}}}}
    {\idxappend{\idxshape{\sequence{\nat_s}} \; \EraseT{\type_B}}}}
\begin{alltt}
\(\expr = \)(unbox (\(\unboxvars\) \(\unboxval\)) \(\unboxbody\))\(\sp{\type\sb{R}}\)\\
\(\mapsto\sb{\mathit{unbox}}\)\\
\(\expr' = \unboxresult\)
\end{alltt}
Then relying on our earlier result that erasure commutes with substitution,
we can step the erased version of $\expr$ to get the erased version of $\expr'$:
\begin{alltt}
\(\EraseE{\expr} = \) (unbox (\(\unboxvars\) \(\erasedunboxval\)) \(\EraseE{\expr\sb{b}}\) \(\EraseT{\type\sb{B}}\))\\
\(\mapsto\sb{\mathit{unbox}} \erasedunboxresult\)\\
\( = \rewrittenunboxresult = \EraseE{\expr'}\)
\end{alltt}

\paragraph{Mis-applied primitive operator:}
The remaining case is that $\expr \not\mapsto$, but $\expr$ is not a value.
Since $\expr$ is well-typed, Lemma \ref{Progress} (Progress) implies that
$\expr$ must be a mis-application of a primitive operator.
That is, $\expr$ has the form
$\ctxt\left[\misapp\right]$,
where $\Sigref{\primop}$ is
{\tt (-> ($\type_I$ $\sequence{}$) $\type_O$)},
and the types of $\sequence{\val}$ are $\sequence{\type_i}$.
Using Lemma \ref{ErasureInContext} (Erasure in context),
$\EraseE{\expr}$ is
$\EraseC{\ctxt}[${\tt ($\primop$ ($\EraseE{\val}$ $\EraseT{\type_I}$) $\sequence{}$ $\EraseT{\type_r}$)}$]$.
Since each $\val$ has shape $\EraseT{\type_I}$,
we still have function application with a scalar frame,
and the values given as arguments are still out-of-domain for $\primop$.
\end{sproof}

Since we have a deterministic operational semantics for
both explicitly typed Remora and type-erased Remora,
the lockstep lemma also works in reverse.
If an erased term takes an evaluation step,
its preimage cannot be a value form or stuck state.
The preimage must therefore step to some result expression,
which itself erases to the same result.
Similarly, a value form or stuck state in erased Remora
cannot have a preimage which takes an evaluation step.

\begin{corollary}[Reverse lockstep]
If $\EraseE{\expr} \mapsto \EraseE{\expr'}$,
then for any $\expr''$ such that $\expr \mapsto \expr''$, we have $\expr' \EraseEqv \expr''$,
and at least one such $\expr''$ exists.
If $\EraseE{\expr} \not\mapsto$, then $\expr \not\mapsto$.
\end{corollary}

Recall our relation $\EraseEqv$ on the set of machine states $S = \wt{\Expr} \uplus \erased{\wt{\Expr}}$,
where $\wt{\Expr} = \{e \in \Expr | \typeof{\emptyEnv}{\emptyEnv}{\emptyEnv}{\expr}{\type}\}$,
\ie the set of well-typed explicitly-typed terms,
and $\erased{\wt{\Expr}}$ is the image of $\wt{\Expr}$ under type erasure.
$\EraseEqv$ is the equivalence closure of the relation given by the erasure function $\EraseE{\cdot}$.
That is, $\EraseEqv$ is the least relation which
relates two states $s$ and $w$ iff any of the following hold:
\begin{enumerate}
\item{$s \in \wt{Expr}$ and $\EraseE{s} = w$ (erasure proper)}
\item{$s =_\alpha w$ (reflexivity)}
\item{$w \EraseEqv s$ (symmetry)}
\item{$s \EraseEqv s'$ and $s' \EraseEqv w$ (transitivity)}
\end{enumerate}
Expanding the erasure relation based on $\EraseE{\cdot}$
to include both symmetry and transitivity
relates any two explicitly typed expressions
which produce $\alpha$-equivalent erased terms.
A $\EraseEqv$ equivalence class consists of a single erased Remora expression
and all of its preimages.
There can be only one erased Remora expression because
type erasure is a well-defined function
(\ie, no single explicitly typed expression
can erase to multiple different results).
Formally, every $\EraseEqv$ equivalence class must have the form
$$\{\eexpr\} \uplus \{\expr \in \Expr \, | \, \EraseE{\expr} = \eexpr\}$$
\begin{theorem}
  \label{Bisimulation}
  $\EraseEqv$ is a bisimulation.
  That is, for any states $s, w \in S$ if $s \EraseEqv w$,
  either $(s \mapsto u \wedge w \mapsto v \wedge u \EraseEqv v)$
  or $(s \not\mapsto \wedge w \not\mapsto)$.
\end{theorem}
\begin{sproof}[\BODY]
  There are four cases to consider,
  depending on which of $\Expr$ or $\erased{\Expr}$ each related term is drawn from,
  but we can merge the two cases where $s$ and $w$ are drawn from different languages.

  \paragraph{$s \in \erased{\Expr}$ and $w \in \Expr$, or vice versa:}
  Then $s$ is the sole type-erased expression in its equivalence class,
  and $\EraseE{w} = s$
  (or vice versa).
  Our proof obligation is exactly the lockstep lemma (Lemma \ref{Lockstep}).

  \paragraph{$s, w \in \erased{\Expr}$:}
  Since each equivalence class contains only one type-erased expression, $s = w$.
  They must therefore have the same reduction behavior.

  \paragraph{$s, w \in \Expr$:}
  If $s \not\mapsto$, then the lockstep lemma implies $\EraseE{s} = \EraseE{w} \not\mapsto$.
  Then by reverse lockstep, $w \not\mapsto$ as well.
  On the other hand, if $s \mapsto s'$, then $\EraseE{s} = \EraseE{w} \mapsto \EraseE{s'}$.
  Lockstep implies $\EraseE{w} \mapsto \EraseE{w'}$.
  Since Erased Remora has deterministic operational semantics,
  $\EraseE{s'} = \EraseE{w'}$
  (they are both the result of taking an evaluation step from the same expression).
  Therefore, all of their preimages, including $s'$ and $w'$ are erasure-equivalent,
  \ie, $s' \EraseEqv w'$.
\end{sproof}

% Now that we have established the same ``behavior'' during evaluation,
% we add one final note showing that the end results are the same
% (\ie, critical details found in an explicit result term
% have not been lost in translation).
% \begin{lemma}[Uniqueness of first-order preimages]
%   \label{FirstOrderPreimage}
%   For any $\val$, $\val'$ and first-order type $\type$
%   where $\typeof{\emptyEnv}{\emptyEnv}{\emptyEnv}{\val}{\type}$
%   and $\typeof{\emptyEnv}{\emptyEnv}{\emptyEnv}{\val'}{\type}$,
%   $\EraseE{\val} = \EraseE{\val'} \implies \val = \val'$.
% \end{lemma}
% \begin{sproof}
%   \pfstub
% \end{sproof}
